\chapter{Partial Waves, Dispersion Relations, and High-Energy Bounds}
\label{ch:volume-ii-07}

\section{Angular Analyticity}

Work in \(D=4\) and consider scattering of the lightest stable identical
scalar particle of mass \(m\).  The connected invariant amplitude is denoted
by
\[
  \mathcal A(s,t),
  \qquad
  s+t+u=4m^2 .
\]
For a fixed physical value \(s>4m^2\), introduce
\[
  x=\cos\theta
  =
  1+\frac{2t}{s-4m^2}.
\]
The physical \(s\)-channel scattering angles give the interval
\[
  -1\le x\le1 .
\]
At the \(t\)-channel two-particle threshold \(t=4m^2\),
\[
  x=x_t(s)\equiv\frac{s+4m^2}{s-4m^2}>1,
\]
and the \(t\)-channel cut maps to \(x\in[x_t,\infty)\).  The \(u\)-channel
threshold gives the symmetric cut \(x\in(-\infty,-x_t]\).  If a crossed
channel contains a stable bound state of mass \(M_B<2m\), its pole appears
at
\[
  x_B=1+\frac{2M_B^2}{s-4m^2}
\]
or at the corresponding point on the left.  Let \(t_0>0\) be the smallest
positive value of \(t\) at which such a crossed-channel singularity occurs.
For large \(s\), the nearest right singularity is
\[
  x_0(s)=1+\frac{2t_0}{s-4m^2}.
\]

\begin{center}
\begin{tikzpicture}[>=Stealth,scale=0.86,every node/.style={font=\small}]
  \draw[->,line width=0.75pt] (-4.8,0)--(5.05,0)
    node[below] {\(\operatorname{Re}x\)};
  \draw[->,line width=0.75pt] (0,-2.15)--(0,2.35)
    node[left] {\(\operatorname{Im}x\)};
  \draw[blue!70!black,line width=1.4pt] (-1,0)--(1,0);
  \fill[blue!70!black] (-1,0) circle (1.8pt);
  \fill[blue!70!black] (1,0) circle (1.8pt);
  \node[blue!70!black,below] at (-1,-0.18) {\(-1\)};
  \node[blue!70!black,below] at (1,-0.18) {\(1\)};
  \node[blue!70!black,fill=white,inner sep=1pt] at (0,0.33)
    {physical interval};

  \draw[red!70!black,decorate,
    decoration={snake,amplitude=0.55mm,segment length=3.2mm},line width=1.0pt]
    (2.25,0)--(4.6,0);
  \draw[red!70!black,decorate,
    decoration={snake,amplitude=0.55mm,segment length=3.2mm},line width=1.0pt]
    (-4.6,0)--(-2.25,0);
  \node[red!70!black,below] at (2.25,-0.24) {\(x_t\)};
  \node[red!70!black,below] at (-2.25,-0.24) {\(-x_t\)};
  \node[red!70!black,align=center] at (3.65,-0.68) {\(t\)-channel\\cut};
  \node[red!70!black,align=center] at (-3.65,0.68) {\(u\)-channel\\cut};

  \fill[orange!85!black] (1.62,0) circle (2.2pt);
  \draw[orange!85!black,->,line width=0.85pt] (2.35,1.28)--(1.69,0.12);
  \node[orange!85!black,align=left,anchor=west] at (2.42,1.28)
    {possible crossed\\bound-state pole};

  \draw[purple!75!black,dashed,line width=1.1pt] (0,0) ellipse (1.9 and 1.25);
  \draw[green!50!black,dashed,line width=0.95pt] (0,0) ellipse (1.35 and 0.72);
  \node[purple!75!black,align=center] at (-2.25,1.55)
    {largest allowed\\ellipse \(\gamma_s\)};
  \node[green!50!black,align=center] at (-1.55,-0.92)
    {smaller ellipse};
  \node[align=center] at (0,-1.75)
    {ellipses have foci at \(x=\pm1\)};
\end{tikzpicture}
\end{center}

The statement used below is the following angular analyticity hypothesis:
for each such \(s\), \(\mathcal A(s,x)\) is analytic in a domain containing
the physical interval and bounded on every compact subset of that domain.
The largest ellipse with foci at \(\pm1\) lying in this analytic domain is
called the Lehmann ellipse at energy \(s\).

\section{Expansion Inside the Lehmann Ellipse}

Let \(\gamma_s\) be a counterclockwise contour inside the angular analytic
domain and enclosing \([-1,1]\).  For \(x\) inside \(\gamma_s\),
\[
  \mathcal A(s,x)
  =
  \oint_{\gamma_s}\frac{\dd z}{2\pi\ii}\,
  \frac{\mathcal A(s,z)}{z-x}.
\]
The kernel has the Legendre expansion
\[
  \frac{1}{z-x}
  =
  \sum_{\ell=0}^{\infty}(2\ell+1)P_\ell(x)Q_\ell(z),
\]
where
\[
  Q_\ell(z)
  =
  \frac12\int_{-1}^{1}\dd y\,\frac{P_\ell(y)}{z-y}.
\]
This expansion converges when
\[
  \left|z+\sqrt{z^2-1}\right|
  >
  \left|x+\sqrt{x^2-1}\right|.
\]
Indeed, at large \(\ell\),
\[
  P_\ell(x)\sim
  \left(x+\sqrt{x^2-1}\right)^\ell,
  \qquad
  Q_\ell(z)\sim
  \left(z+\sqrt{z^2-1}\right)^{-\ell},
\]
and the level sets of
\[
  \left|z+\sqrt{z^2-1}\right|
\]
are ellipses with foci at \(\pm1\).  Thus angular analyticity implies
convergence of the partial-wave expansion throughout every smaller ellipse
inside the Lehmann ellipse.

\section{Partial-Wave Normalization}

For identical scalar particles only even partial waves occur.  Let
\(\mathcal M_{\rm ord}(s,x)\) denote the ordered two-particle amplitude whose
partial-wave expansion is
\[
  \mathcal M_{\rm ord}(s,x)
  =
  16\pi
  \sum_{\ell=0}^{\infty}
  (2\ell+1)a_\ell(s)P_\ell(x).
\]
The amplitude used in this chapter is the Bose-symmetrized amplitude on the
space of identical two-particle states:
\[
  \mathcal A(s,x)
  =
  \mathcal M_{\rm ord}(s,x)+\mathcal M_{\rm ord}(s,-x)
  =
  32\pi
  \sum_{\substack{\ell=0\\ \ell\ {\rm even}}}^{\infty}
  (2\ell+1)a_\ell(s)P_\ell(x).
\]
Using
\[
  S_\ell(s)=1+2\ii\,\rho(s)a_\ell(s),
  \qquad
  \rho(s)=\sqrt{1-\frac{4m^2}{s}},
\]
this becomes the convention
\[
  \mathcal A(s,x)
  =
  -16\pi\ii\,\rho(s)^{-1}
  \sum_{\substack{\ell=0\\ \ell\ {\rm even}}}^{\infty}
  (2\ell+1)\bigl(S_\ell(s)-1\bigr)P_\ell(x).
\]
Here \(S_\ell(s)\) is the partial-wave eigenvalue of the \(S\)-operator in the
two-particle sector.  In an elastic interval \(|S_\ell|=1\).  With inelastic
channels present, unitarity gives
\[
  |S_\ell(s)|\le1 .
\]
For real \(x\) inside the Lehmann ellipse, real analyticity gives
\[
  \operatorname{Im}\mathcal A(s,x)
  =
  16\pi\,\rho(s)^{-1}
  \sum_{\substack{\ell=0\\ \ell\ {\rm even}}}^{\infty}
  (2\ell+1)P_\ell(x)\,w_\ell(s),
\]
where
\[
  w_\ell(s)=1-\operatorname{Re}S_\ell(s).
\]
The inequality \(|S_\ell|\le1\) implies the positive bound
\[
  0\le w_\ell(s)\le2 .
\]

The elastic and inelastic cross sections are
\[
  \sigma_{\rm el}
  =
  \frac{8\pi}{s-4m^2}
  \sum_{\substack{\ell=0\\ \ell\ {\rm even}}}^{\infty}
  (2\ell+1)|S_\ell(s)-1|^2
\]
and
\[
  \sigma_{\rm inel}
  =
  \frac{8\pi}{s-4m^2}
  \sum_{\substack{\ell=0\\ \ell\ {\rm even}}}^{\infty}
  (2\ell+1)\bigl(1-|S_\ell(s)|^2\bigr).
\]
Adding them gives
\[
  \sigma_{\rm tot}
  =
  \frac{16\pi}{s-4m^2}
  \sum_{\substack{\ell=0\\ \ell\ {\rm even}}}^{\infty}
  (2\ell+1)w_\ell(s).
\]
Equivalently,
\[
  \sigma_{\rm tot}
  =
  \frac{1}{\sqrt{s(s-4m^2)}}\,
  \operatorname{Im}\mathcal A(s,1).
\]
This is the optical theorem in the present normalization.

\section{The Froissart--Martin Mechanism}

Define the positive total weight
\[
  \bar L^2(s)
  =
  \sum_{\substack{\ell=0\\ \ell\ {\rm even}}}^{\infty}
  (2\ell+1)w_\ell(s)
  =
  \frac{s-4m^2}{16\pi}\,\sigma_{\rm tot}(s).
\]
For real \(x>1\), \(P_\ell(x)>0\) and increases with \(\ell\).  Among all
weights satisfying \(0\le w_\ell\le2\) and fixed total \(\bar L^2\), the
smallest value of
\[
  \sum_{\ell\ {\rm even}}(2\ell+1)P_\ell(x)w_\ell
\]
is obtained, up to endpoint rounding, by filling the lowest allowed angular
momenta to height \(2\).  Denote the corresponding box profile by \(b_\ell\).
At high energy the endpoint \(L\) of the box obeys
\[
  L^2\simeq \bar L^2 .
\]

\begin{center}
\begin{tikzpicture}[>=Stealth,scale=0.86,every node/.style={font=\small}]
  \begin{scope}
    \draw[->,line width=0.75pt] (-2.9,0)--(3.15,0)
      node[below] {\(\operatorname{Re}x\)};
    \draw[->,line width=0.75pt] (0,-1.35)--(0,1.58)
      node[left] {\(\operatorname{Im}x\)};
    \draw[purple!75!black,dashed,line width=1.0pt] (0,0) ellipse (2.05 and 1.05);
    \draw[blue!70!black,line width=1.2pt] (-1,0)--(1,0);
    \fill[blue!70!black] (-1,0) circle (1.7pt);
    \fill[blue!70!black] (1,0) circle (1.7pt);
    \fill[green!50!black] (1.43,0) circle (2.1pt);
    \node[green!50!black,align=center] at (1.43,0.55)
      {\(x>1\)\\inside \(\gamma_s\)};
    \node[purple!75!black] at (-0.2,1.28) {\(\gamma_s\)};
  \end{scope}
  \begin{scope}[xshift=6.7cm]
    \draw[->,line width=0.75pt] (-0.25,0)--(4.35,0) node[below] {\(\ell\)};
    \draw[->,line width=0.75pt] (0,-0.12)--(0,2.1);
    \node[anchor=east] at (-0.1,1.85) {\(w_\ell\)};
    \draw[purple!75!black,line width=1.0pt]
      (0,1.55)--(0.45,1.55)--(0.45,1.1)--(0.95,1.1)--(0.95,1.38)
      --(1.45,1.38)--(1.45,0.88)--(2.0,0.88)--(2.0,1.18)--(2.6,1.18)
      --(2.6,0.45)--(3.2,0.45)--(3.2,0.22)--(3.85,0.22);
    \draw[blue!70!black,line width=1.1pt] (0,1.45)--(2.75,1.45)--(2.75,0);
    \draw[red!70!black,dashed,line width=0.85pt] (2.75,0)--(2.75,1.72);
    \node[blue!70!black,fill=white,inner sep=1pt] at (1.33,0.68)
      {box \(b_\ell\)};
    \node[red!70!black,below] at (2.75,-0.06) {\(L\)};
    \node[purple!75!black,align=center] at (3.0,1.9)
      {allowed\\weights};
  \end{scope}
\end{tikzpicture}
\end{center}

For \(x>1\), the integral representation
\[
  P_\ell(x)=\frac{1}{\pi}\int_0^\pi
  \left(x+\cos\alpha\,\sqrt{x^2-1}\right)^\ell\dd\alpha
\]
implies the following lower bound: for every fixed small \(\delta>0\) there
is a constant \(C_\delta>0\) such that
\[
  P_\ell(x)\ge
  C_\delta
  \left(1+(1-\delta)\sqrt{x^2-1}\right)^\ell
  \equiv C_\delta y^\ell .
\]
Using the box profile up to \(L\),
\[
  \operatorname{Im}\mathcal A(s,x)
  \ge
  C\,\rho(s)^{-1}
  \sum_{\substack{0\le\ell\le L\\ \ell\ {\rm even}}}
  (2\ell+1)y^\ell
  \ge
  C'\,\rho(s)^{-1}L\,\frac{y^L}{y-1},
\]
with constants independent of \(s\) at sufficiently large energy.

Now choose \(x=x_0(s)=1+2t_0/(s-4m^2)\), the edge of the Lehmann ellipse set by
the nearest positive-\(t\) singularity.  At large \(s\),
\[
  y-1=(1-\delta)\sqrt{x_0^2-1}+O(s^{-1})
  =
  2(1-\delta)\sqrt{\frac{t_0}{s}}+O(s^{-1}),
\]
and
\[
  L\simeq\sqrt{\frac{s\,\sigma_{\rm tot}}{16\pi}} .
\]
Therefore the exponential factor obeys
\[
  y^L
  \simeq
  \exp\!\left[
    (1-\delta)
    \sqrt{\frac{t_0\,\sigma_{\rm tot}}{4\pi}}
  \right].
\]
If, at the ellipse edge,
\[
  |\mathcal A(s,x_0(s))|\le C_Ns^N
  \qquad(s\to\infty),
\]
then the exponential lower bound is compatible with this polynomial upper
bound only if
\[
  \sigma_{\rm tot}(s)
  \le
  \frac{(N-1)^2}{(1-\delta)^2}\,
  \frac{4\pi}{t_0}
  \left(\log\frac{s}{M_0^2}\right)^2
\]
for sufficiently large \(s\).  The scale \(M_0\) changes only subleading
terms in the logarithm.  This is the Froissart--Martin mechanism: angular
analyticity limits the number of strongly contributing partial waves, and
unitarity bounds the weight carried by each partial wave.

\section{Polynomial Boundedness}

The polynomial boundedness assumption has a causal interpretation.  A
one-dimensional model already exhibits the relevant analytic structure.  Let
\[
  f_{\rm out}(t)=\int\dd t'\,K(t-t')f_{\rm in}(t'),
  \qquad
  \widetilde K(\omega)=\int\dd t\,\ee^{\ii\omega t}K(t).
\]
If the kernel has retarded support,
\[
  K(t)=0\qquad(t<0),
\]
then \(\widetilde K(\omega)\) is analytic for \(\operatorname{Im}\omega>0\).
Moreover a growth like \(\exp(c\,\operatorname{Im}\omega)\) in the upper half
plane would correspond to support at negative time.  Relativistic scattering
does not reduce to this one-dimensional model, but the model explains the
form of the hypothesis: locality and causal propagation are reflected in
analyticity and in the absence of exponential growth in complex energy
directions.

For the invariant amplitude, fix real \(t<0\) and continue \(s\) to the
complex plane.  The displayed dispersion formula records the cut terms in a
kinematic situation where bound-state pole terms and anomalous first-sheet
terms are absent; when such terms are present, their residues are added
explicitly.  Assume that for some integer \(N\),
\[
  \lim_{|s|\to\infty}
  \left|\frac{\mathcal A(s,t)}{s^N}\right|=0
\]
in the domain used to deform the contour.  Choose subtraction points
\[
  s_1,\ldots,s_N
\]
away from the cuts and from the evaluation point \(s\).  Then
\[
  \frac{\mathcal A(s',t)}{\prod_{j=1}^N(s'-s_j)}
\]
has sufficient decay for the large circle to be removed in Cauchy's theorem.
Equivalently,
\[
  \mathcal A(s,t)
  =
  \oint_{\gamma}
  \frac{\dd s'}{2\pi\ii}\,
  \frac{\mathcal A(s',t)}{s'-s}
  \prod_{j=1}^N\frac{s-s_j}{s'-s_j},
\]
where \(\gamma\) encloses \(s\) and the subtraction points before deformation.

\begin{center}
\begin{tikzpicture}[>=Stealth,scale=0.86,every node/.style={font=\small}]
  \draw[->,line width=0.75pt] (-4.35,0)--(4.55,0)
    node[below] {\(\operatorname{Re}s'\)};
  \draw[->,line width=0.75pt] (0,-2.15)--(0,2.35)
    node[left] {\(\operatorname{Im}s'\)};
  \draw[red!70!black,decorate,
    decoration={snake,amplitude=0.55mm,segment length=3.1mm},line width=1.0pt]
    (1.35,0)--(4.1,0);
  \draw[red!70!black,decorate,
    decoration={snake,amplitude=0.55mm,segment length=3.1mm},line width=1.0pt]
    (-4.1,0)--(-0.85,0);
  \node[red!70!black,below] at (1.35,-0.08) {\(4m^2\)};
  \node[red!70!black,below] at (-0.85,-0.08) {\(-t\)};
  \node[red!70!black,align=center] at (3.05,-0.55) {\(s\)-channel\\cut};
  \node[red!70!black,align=center] at (-2.7,0.55) {crossed\\cut};

  \draw[purple!75!black,line width=0.95pt,->]
    (0,0) ++(1.55,0) arc (0:335:1.55);
  \node[purple!75!black] at (1.45,1.5) {\(\gamma\)};
  \fill[black] (0.42,0.55) circle (1.8pt);
  \node[anchor=west] at (0.52,0.55) {\(s\)};
  \foreach \x/\lab in {-0.25/s_1,0.92/s_2,1.88/s_N} {
    \draw[green!50!black,line width=0.85pt] (\x,0) circle (0.16);
    \node[green!50!black,above] at (\x,0.22) {\(\lab\)};
  }
  \draw[blue!70!black,->,line width=0.9pt] (1.4,0.32)--(3.75,0.32);
  \draw[blue!70!black,->,line width=0.9pt] (-0.85,-0.32)--(-3.75,-0.32);
  \node[blue!70!black,align=center,fill=white,inner sep=1.5pt] at (0,-1.88)
    {deformed contour:\\cuts plus subtraction residues};
\end{tikzpicture}
\end{center}

Let
\[
  \operatorname{disc}\mathcal A(s,t)
  =
  \mathcal A(s+\ii0,t)-\mathcal A(s-\ii0,t).
\]
After deforming \(\gamma\) onto the right and left cuts, one obtains
\[
  \mathcal A(s,t)
  =
  P_{N-1}(s;t)
  +
  \frac{1}{2\pi\ii}
  \int_{C_s\cup C_u}
  \dd s'\,
  \operatorname{disc}\mathcal A(s',t)\,
  \frac{1}{s'-s}
  \prod_{j=1}^N\frac{s-s_j}{s'-s_j}.
\]
Here \(C_s=[4m^2,\infty)\) is the right cut, \(C_u=(-\infty,-t]\) is the
crossed cut in the \(s'\)-plane, and \(P_{N-1}(s;t)\) is the polynomial of
degree at most \(N-1\) produced by the small circles around the subtraction
points.  On a real cut where real analyticity holds,
\[
  \operatorname{disc}\mathcal A(s,t)
  =
  2\ii\,\operatorname{Im}\mathcal A(s+\ii0,t).
\]
The formula is an \(N\)-subtracted fixed-\(t\) dispersion relation.

\section{Two Subtractions}

The previous derivation used a complex-plane polynomial bound.  A useful
refinement is that control of the absorptive part on the real axis determines
how many subtractions are needed.  Suppose
\[
  |\operatorname{Im}\mathcal A(s,t)|\le C s^{N-\epsilon}
  \qquad(s\to+\infty)
\]
for some \(\epsilon>0\).  If a dispersion formula with \(N'\ge N+1\)
subtractions were genuinely required, its polynomial part would generically
give
\[
  \mathcal A(s,t)\sim C'(t)s^{N'-1}.
\]
Unitarity at \(t=0\) gives
\[
  \operatorname{Im}\mathcal A(s,0)
  =
  \sqrt{s(s-4m^2)}\,\sigma_{\rm tot}(s)
  \ge
  \sqrt{s(s-4m^2)}\,\sigma_{\rm el}(s).
\]
The elastic cross section contains the positive integral of
\(|\mathcal A(s,t)|^2\) over the physical angular interval.  A leading growth
like \(s^{N'-1}\) would therefore force at least
\[
  \operatorname{Im}\mathcal A(s,0)
  \ge c\,s^{2N'-2}
\]
for some positive constant \(c\) along a suitable high-energy sequence,
contradicting the assumed real-axis bound.  Thus the absorptive part controls
the subtraction number.

The Froissart estimate gives, at forward scattering,
\[
  \operatorname{Im}\mathcal A(s,0)
  \le
  C\,s(\log s)^2 .
\]
For fixed real \(t<0\) and sufficiently large \(s\), the corresponding angular
variable satisfies \(-1<x<1\).  Since \(P_\ell(x)\le P_\ell(1)=1\) in this
range,
\[
  \operatorname{Im}\mathcal A(s,t)
  \le
  \operatorname{Im}\mathcal A(s,0)
\]
for the positive partial-wave sum above.  Consequently
\[
  \lim_{s\to\infty}
  \left|\frac{\mathcal A(s,t)}{s^2}\right|=0,
  \qquad t<0,
\]
so two subtractions are sufficient in the massive scalar setting.

There is a further extension from real \(t<0\) to complex \(t\) in the disk
set by the nearest \(t\)-plane singularity.  The mechanism is positivity of
the Taylor coefficients of the absorptive part at \(t=0\):
\[
  \left.
  \frac{\partial^n}{\partial t^n}
  \operatorname{Im}\mathcal A(s,t)
  \right|_{t=0}
  \ge0 .
\]
This positivity allows the Taylor expansion in \(t\) to commute with the
dispersive integral inside its convergence disk.  With two subtractions the
Froissart--Martin bound takes the standard four-dimensional form
\[
  \sigma_{\rm tot}(s)
  \le
  \frac{1}{(1-\delta)^2}\,
  \frac{4\pi}{t_0}
  \left(\log\frac{s}{M_0^2}\right)^2
\]
for sufficiently large \(s\).  In \(D\) spacetime dimensions the corresponding
geometric power is
\[
  \sigma_{\rm tot}(s)\le C\,(\log s)^{D-2}.
\]

\section{Scope of the Bound}

The bound is a theorem about amplitudes satisfying the stated massive
analyticity, unitarity, crossing, and boundedness hypotheses.  The mass gap
is part of the statement.  When massless particles are present, the total
cross section may diverge because arbitrarily small scattering angles or
collinear emissions are experimentally unresolved.  The corresponding
observable is then an inclusive or angularly cut cross section, not the total
cross section used above.

Gravity also lies outside the same massive-QFT hypothesis.  In \(D\) spacetime
dimensions, a black hole of energy \(E\) has Schwarzschild radius
\[
  r_H\sim (G_NE)^{1/(D-3)},
\]
so the effective interaction size grows as a power of energy.  This is a
different high-energy problem, not a modification of the massive
Froissart--Martin argument.
